\documentclass[acmtog]{acmart}
\citestyle{acmauthoryear}
\usepackage{graphicx}
\usepackage{overpic}
\usepackage{caption}
\usepackage{subcaption}
\usepackage{float}
\usepackage{xcolor}
\usepackage{tabularx}
\usepackage{algorithm}
\usepackage{algpseudocode}
%%
%% \BibTeX command to typeset BibTeX logo in the docs
% \AtBeginDocument{%
%   \providecommand\BibTeX{{%
%     Bib\TeX}}}

\setcopyright{cc}
\setcctype{by}
\acmJournal{TOG}
\acmYear{2026} \acmVolume{45} \acmNumber{6} \acmArticle{184}
\acmMonth{12} \acmDOI{10.1145/3842537}




\begin{document}


\title[Hypergraph Rotation Systems]{Hypergraph Rotation Systems: \\A Unified Framework for Linked Knot Structures}


\author{Tolga Talha Yıldız}
\affiliation{%
  \institution{Texas A\&M University}
  \city{College Station}
  \state{Texas}
  \country{USA}
}
\email{tolgayildiz@tamu.edu}
\orcid{0000-0002-4834-9389}

\author{Uğur Önal}
\affiliation{%
  \institution{Texas A\&M University}
  \city{College Station}
  \state{Texas}
  \country{USA}
}
\email{uguronal@tamu.edu}
\orcid{0000-0001-8347-3038}

\author{Alison Martin}
\affiliation{%
  \institution{Weaver and Independent Researcher}
  \city{Livorno}
  \country{Italy}
}
\email{alisonmartin57@gmail.com}
\orcid{0000-0003-0207-3898}

\author{Vinayak R. Krishnamurthy}
\affiliation{%
  \institution{Texas A\&M University}
  \city{College Station}
  \state{Texas}
  \country{USA}
}
\email{vinayak@tamu.edu}
\orcid{0000-0001-6434-2577}

\author{Ergun Akleman}
\affiliation{%
  \institution{Texas A\&M University}
  \city{College Station}
  \state{Texas}
  \country{USA}
}
\authornote{Corresponding Author}
\email{ergun@tamu.edu}
\orcid{0000-0003-3618-4166}

\renewcommand{\shortauthors}{Yildiz et al.}


\begin{abstract}

We present a unified framework for the design of linked knot (LK) structures including links, knots, weaves, knits, and chainmail structures, extending them to articulated and compliant (spring-like) assemblies. The core of our approach is a combinatorial representation based on \emph{Hypergraph Rotation Systems (H-GRS)}, in which topology is defined by the cyclic ordering of incidences between faces and edges. This formulation removes the need for manifold assumptions or global orientation, naturally supporting non-manifold connectivity.

Building on this representation, we introduce integer-valued twist assignments on edges as local operators that deterministically control global cycle structure, linking, and connectivity. These operators enable the systematic construction of both unicursal (single-cycle) and multi-component configurations, which can exhibit different forms of compliant and articulated behavior.

Unlike prior approaches based on manifold constraints or binary twisting, our framework supports multi-valued twisting and higher-order connectivity, enabling selective linking, partial merging, and articulated couplings not achievable in classical textile constructions. 

We demonstrate the framework across a wide range of structures and validate it through physical prototypes fabricated via 3D printing and paper-strip construction on a scaffold. By unifying combinatorial topology with the exploration of mechanical behavior, H-GRS provides a minimal and extensible foundation for the topology-driven design of linked filament structures, with applications in computational fabrication, architectured materials, and bio-inspired systems.

\end{abstract}

%%
%% The code below is generated by the tool at http://dl.acm.org/ccs.cfm.
%% Please copy and paste the code instead of the example below.
%%
\begin{CCSXML}
<ccs2012>
   <concept>
       <concept_id>10010147.10010371.10010396.10010398</concept_id>
       <concept_desc>Computing methodologies~Mesh geometry models</concept_desc>
       <concept_significance>500</concept_significance>
       </concept>
   <concept>
       <concept_id>10010147.10010371.10010396.10010399</concept_id>
       <concept_desc>Computing methodologies~Parametric curve and surface models</concept_desc>
       <concept_significance>500</concept_significance>
       </concept>
 </ccs2012>
\end{CCSXML}
\ccsdesc[500]{Computing methodologies~Mesh geometry models}
\ccsdesc[500]{Computing methodologies~Parametric curve and surface models}

\keywords{hypergraph rotation systems, linked knots, linked structures,
knots and links, topological modeling, computational design}




\begin{teaserfigure}
    \centering
    \begin{subfigure}[t]{0.16\textwidth}
        \includegraphics[width=\textwidth]{images/teaser/cube/cube_base.jpg}
        \includegraphics[width=\textwidth]{images/teaser/cube/2man_chainmail.jpg}
        \caption{A link.}
        \label{teaser0}
    \end{subfigure}
    \hfill
        \begin{subfigure}[t]{0.16\textwidth}
        \includegraphics[width=\textwidth]{images/teaser/cube/cube_base.jpg}
        \includegraphics[width=\textwidth]{images/teaser/cube/2man_knot.jpg}
        \caption{A knot.}
        \label{teaser1}
    \end{subfigure}
    \hfill
    \begin{subfigure}[t]{0.16\textwidth}
        \includegraphics[width=\textwidth]{images/teaser/cube/unfolded_base.jpg}
        \includegraphics[width=\textwidth]{images/teaser/cube/unfolded.jpg}
        \caption{LK structures.}
        \label{teaser2}
    \end{subfigure}
    \hfill
    \begin{subfigure}[t]{0.16\textwidth}
        \includegraphics[width=\textwidth]{images/teaser/cube/lid_base.jpg}
        \includegraphics[width=\textwidth]{images/teaser/cube/lid.jpg}
        \caption{LK structures.}
        \label{teaser3}
    \end{subfigure}
    \hfill
    \begin{subfigure}[t]{0.16\textwidth}
        \includegraphics[width=\textwidth]{images/teaser/cube/piano1_base.jpg}
        \includegraphics[width=\textwidth]{images/teaser/cube/piano1.jpg}
        \caption{LK structures.}
        \label{teaser4}
    \end{subfigure}
    \hfill
    \begin{subfigure}[t]{0.16\textwidth}     
    \includegraphics[width=\textwidth]{images/teaser/cube/cp_base.jpg}
    \includegraphics[width=\textwidth]{images/teaser/cube/cp.jpg}
        \caption{Periodic LKs.}
        \label{teaser5}
    \end{subfigure}
    \hfill
  \caption{
  Local cyclic face orderings and integer-labeled edges provide a minimal representation for designing linked structures with different topology and motion. Starting from simple scaffolds (top), the resulting Linked Knot (LK) geometries (bottom) exhibit a wide range of behaviors, including chainmail links (a), unicursal knots (b), articulated structures (c–d), functional mechanisms (e), and periodic configurations (f). These examples illustrate how variations in scaffold connectivity, boundary conditions, periodicity, and local twist assignments generate distinct LK structures within a single H-GRS framework.}
  \Description{Six paired examples show scaffold geometries above and their corresponding linked-knot structures below, including links, knots, articulated mechanisms, and a periodic lattice. Distinct strand colors indicate separate components.}
  \label{fig:teaser}
\end{teaserfigure}
\received{02 May 2026}
\received[revised]{19 August 2026}
\received[accepted]{25 August 2026}

\maketitle

\begin{figure*}[htb!]
    \centering
    \includegraphics[width=0.19\textwidth]{images/2man/input_mesh.jpg}
    \includegraphics[width=0.19\textwidth]{images/2man/n4_tw.jpg}
    \includegraphics[width=0.19\textwidth]{images/2man/n3_tw.jpg}
    \includegraphics[width=0.19\textwidth]{images/2man/n2_tw.jpg}
    \includegraphics[width=0.19\textwidth]{images/2man/n1_tw.jpg}\\
    \parbox[t]{0.19\textwidth}{\centering Input mesh}
    \parbox[t]{0.19\textwidth}{\centering -4 twist}
    \parbox[t]{0.19\textwidth}{\centering -3 twist}
    \parbox[t]{0.19\textwidth}{\centering -2 twist}
    \parbox[t]{0.19\textwidth}{\centering -1 twist}\\
    \includegraphics[width=0.19\textwidth]{images/2man/0_tw.jpg}
    \includegraphics[width=0.19\textwidth]{images/2man/1_tw.jpg}
    \includegraphics[width=0.19\textwidth]{images/2man/2_tw.jpg}
    \includegraphics[width=0.19\textwidth]{images/2man/3_tw.jpg}
    \includegraphics[width=0.19\textwidth]{images/2man/4_tw.jpg}\\
    \parbox[t]{0.19\textwidth}{\centering 0 twist}
    \parbox[t]{0.19\textwidth}{\centering 1 twist}
    \parbox[t]{0.19\textwidth}{\centering 2 twist}
    \parbox[t]{0.19\textwidth}{\centering 3 twist}
    \parbox[t]{0.19\textwidth}{\centering 4 twist}\\
    \caption{This example shows the effect of twist-labels on an edge shared by two faces. Even twists result in disconnected components; non-zero even twists generate interlocked rings. An odd number of twists generates single-cycle curves, resulting in a knotted surface.}
    \Description{A grid shows an input mesh and linked-knot structures generated with twist values from -4 to 4. Even twists produce two separate loops, while odd twists produce a single connected loop; positive and negative twists form mirrored configurations.}
    \label{fig:2_mantwist}
\end{figure*}



\section{Introduction}
\label{sec:intro}

Linked, knotted, and woven structures appear across a remarkable range of scales and disciplines, emerging from physical, biological, and chemical processes alike. In engineered materials, interlaced fiber architectures arise from frictional contact, hierarchical arrangement, and deformation from textiles and ropes \citep{knittel2020modelling,niu2025geometric,singal2024programming} to deployable architectural gridshells \citep{yang2025topologicalwovennodes,liu2025doublelayeredegridshells} and entangled metamaterial trusses exhibiting high stiffness and stretchability \citep{surjadi2025doublenetwork,moestopo2023knots}. Biology offers equally intricate examples: protein knot motifs and folded peptide structures \citep{craik2001cystine}, fibrillous assemblies in butterfly wings \citep{jessop2025butterfly}, and DNA-based molecular constructs ranging from catenated polyhedra \citep{zhang1994dnatroctahedron} to DNA origami nanostructures \citep{glaser2021dnananostructures}. In chemistry, analogous entanglements emerge as synthetic and self-assembled molecular knots \citep{horner2016knot,fielden2017molecular,forgan2011chemical}. Despite the breadth and complexity of these structures, dedicated modeling frameworks remain scarce and highly case-specific\citep{carton2026threedwovendesign} or limited to two dimensional structures\citep{sonia2026constructionperiodic}. This gap motivates a general computational framework for specifying, modeling, and controlling fibrillar 3D structures from mesh inputs.

In this paper, we introduce a combinatorial framework based on \emph{Hypergraph Rotation Systems (H-GRS)} that decouples structural definition from geometric realization (Figure~\ref{fig:teaser}). Our approach instead leverages topological structure directly, enabling simple, explicit, and controllable design of linked knot structures at the level of discrete mesh topology, rather than mimicking physical or chemical mechanisms. In our formulation, edges are treated as primary entities, and faces define higher-order relationships among them, leading naturally to a hypergraph structure. Local rotation information encodes adjacency and cyclic ordering, allowing global connectivity to emerge through purely combinatorial rules.

Building on this perspective, we demonstrate that linked knot (LK) structures can be constructed through simple local operations on the underlying H-GRS representation. These operations act directly on the combinatorial structure, avoiding the need to simulate physical or chemical processes that typically generate entanglement \citep{moestopo2023knots}. As a result, our approach enables explicit, controllable, and scalable design of complex linked structures.

Furthermore, by separating combinatorial structure from embedding, our framework provides a unified way to describe and generate a wide variety of linked configurations, including those inspired by periodic entanglements, woven systems, and molecular assemblies \citep{evans2013periodiceentanglement1,evans2013periodiceentanglement2,yildiz2025woven,yildiz2025linked}. 
This provides a flexible design space in which linked
configurations can be specified and controlled directly at the combinatorial
level (see Table~\ref{table1}).

While one-dimensional knots become trivial in higher dimensions~\cite{freedman1990topology}, knotted surfaces in four dimensions remain nontrivial~\cite{kamada2002surfaces}. Our construction can be interpreted in this context: we operate on surface-like combinatorial structures whose boundaries are realized as immersed curves in three-dimensional space, giving rise to linked knot structures

\begin{table*}[t]
\centering

\caption{ Comparison of earlier EGRS-based weaving constructions, Block Meshes, and H-GRS. Edge--face incidence and weave foldness are listed separately since, although related in these constructions, they describe distinct properties.}
\label{tab:egrs-hgrs-comparison}
\small
\setlength{\tabcolsep}{4pt}
\begin{tabularx}{\textwidth}{@{}p{0.18\textwidth}p{0.16\textwidth}p{0.11\textwidth}p{0.12\textwidth}p{0.15\textwidth}X@{}}
\toprule
\textbf{Framework / work} &
\textbf{Edge--face incidence} &
\textbf{Weave foldness} &
\textbf{Twist values used} &
\textbf{Manifold requirement} &
\textbf{Main capability} \\
\midrule
EGRS plain weaving~\cite{akleman2009cyclic,xing2010single} &
Two face-sides per edge &
2-fold &
$0,+1$ &
2-manifold surface &
Plain and single-cycle plain weaving \\

EGRS general weaving~\cite{akleman2011cyclic,akleman2015extended,akleman2020topologically} &
Two face-sides per edge &
2-fold &
$0,+1,-1$ &
2-manifold surface &
Twill and other cyclic weaves \\

Block Meshes~\cite{akleman2015blockmeshes} &
Arbitrary number of cyclically ordered face-sides &
Not explicitly discussed &
Not explicitly discussed &
Requires 3-manifold embedding &
Knots and links implicit as side results of a broader theoretical framework \\\\

\textbf{H-GRS (this work)} &
\textbf{Arbitrary number of cyclically ordered face-sides} &
\textbf{Arbitrary-fold} &
$\boldsymbol{t(e)\in\mathbb{Z}}$ &
\textbf{None} &
\textbf{Linked-knot structures} \\
\bottomrule
\end{tabularx}
\label{table1}
\end{table*}



Our contributions are:
\begin{itemize}
      \item The Hypergraph Rotation System (H-GRS), a combinatorial representation of linked knot connectivity through cyclic edge--face incidences, which removes the manifold and global-orientation assumptions of prior rotation-system methods.
      \item Integer-valued edge twists as local operators on the H-GRS that deterministically control global cycle structure; we show (Theorem~\ref{theo:twist}) that the local cycle contribution at an edge of valence $m$ with twist $k$ is governed by $\gcd(m,k)$.
      \item A unified treatment of chainmail, unicursal knots, articulated and compliant mechanisms, and triply periodic woven and knitted materials within a single framework, including a region-decomposition method that separates local compliance from global articulation.
      \item Physical realization of the framework via 3D printing (single-piece, unfolded, and modular) and paper-strip construction, with fabricated prototypes demonstrating articulation, compliance, and periodic assembly.
  \end{itemize}


\begin{figure*}[htb!]
    \centering    
    \includegraphics[width=0.19\textwidth]{images/3man/input_mesh.jpg}    
    \includegraphics[width=0.19\textwidth]{images/3man/n4_tw.jpg}    
    \includegraphics[width=0.19\textwidth]{images/3man/n3_tw.jpg}    
    \includegraphics[width=0.19\textwidth]{images/3man/n2_tw.jpg}    
    \includegraphics[width=0.19\textwidth]{images/3man/n1_tw.jpg}\\
    \parbox[t]{0.19\textwidth}{\centering Input mesh}
    \parbox[t]{0.19\textwidth}{\centering -4 twist}
    \parbox[t]{0.19\textwidth}{\centering -3 twist}
    \parbox[t]{0.19\textwidth}{\centering -2 twist}
    \parbox[t]{0.19\textwidth}{\centering -1 twist}\\    
    \includegraphics[width=0.19\textwidth]{images/3man/0_tw.jpg}    
    \includegraphics[width=0.19\textwidth]{images/3man/1_tw.jpg}    
    \includegraphics[width=0.19\textwidth]{images/3man/2_tw.jpg}   
    \includegraphics[width=0.19\textwidth]{images/3man/3_tw.jpg}   
    \includegraphics[width=0.19\textwidth]{images/3man/4_tw.jpg}\\
    \parbox[t]{0.19\textwidth}{\centering 0 twist}
    \parbox[t]{0.19\textwidth}{\centering 1 twist}
    \parbox[t]{0.19\textwidth}{\centering 2 twist}
    \parbox[t]{0.19\textwidth}{\centering 3 twist}
    \parbox[t]{0.19\textwidth}{\centering 4 twist}\\
    \caption{These examples show the effect of twist-labels on an edge shared by three faces. The twists that are multiples of three result in disconnected but interlocked rings. Any other twist number creates a single connected cycle by combining three original cycles.}
    \Description{A grid shows an input mesh with three faces meeting at one edge and the resulting linked-knot structures for twist values from -4 to 4. Multiples of three produce three separate loops, while other twists produce a single connected loop; opposite signs give mirrored configurations.}
    \label{fig:3_mantwist}
\end{figure*}



\section{Related Work}
\label{sec:related}

Knots and links have a long history in mathematics and topology, with classical work primarily focused on classification and equivalence of embeddings rather than construction \citep{tait1877knots,rolfsen1976knots,adams1994knot,tait1877links,conway1970enumeration,neuwirth1979theory}. While these studies provide deep theoretical insight, they offer limited control for design.

From a constructive perspective, only a few general approaches exist \citep{tait1877knots,akleman2015extended,akleman2020topologically,akleman2015blockmeshes}. These approaches can be interpreted as relying on underlying scaffolds such as graphs, surfaces, or higher-dimensional structures that determine cycles and linking behavior. However, each scaffold introduces limitations in controllability and expressiveness. Prior work on periodic entanglements, volumetric weaving, and linked structures demonstrates rich families of examples but does not provide a unified and extensible design framework \citep{evans2013periodiceentanglement1,evans2013periodiceentanglement2,yildiz2025woven,yildiz2025linked}.

Many practical approaches in engineering and design avoid explicit topological control altogether. Textile-based systems, elastic structures, and discrete assemblies generate interlacing through local rules or physical behavior rather than global structure \citep{borhani2016material,borhani2016emergence,kozlov2013structures,liu2025doublelayeredegridshells,ren2021curvedribbons,verhoeff2011chainlinkfence,roelofs2007rings,yang2025topologicalwovennodes}. Similarly, survey and catalog efforts document linked configurations across domains, including DNA nanostructures and architectural forms, but focus on classification rather than design methodology \citep{glaser2021dnananostructures,kozlov2018knots,gailiunas2022tplinks}.

A key limitation of graph-based constructions, originating from Tait \citep{tait1877knots}, is that an abstract graph does not uniquely determine cycles. Cycles become well-defined only after specifying a rotation system \citep{edmonds1960combinatorial}, and different rotation systems can lead to different knot structures \citep{gross2001topological,akleman1999guaranteeing}. This ambiguity motivates the use of richer representations.

Rotation systems are explicitly encoded in combinatorial surface representations such as winged-edge, half-edge, and quad-edge structures \citep{baumgart1972,mantyla1988,guibas1985,akleman1999guaranteeing}, which define 2-manifold meshes. These representations enable deterministic construction of linked structures by incorporating local twisting information \citep{akleman2009cyclic,akleman2015extended,xing2010single,lipschutz2022linkedknots,sequin2021polyhedral}. 

The main limitation of these prior approaches is an implicit assumption that twisting requires a manifold surface to obtain a consistent orientation. In these approaches, chirality is derived from the inside–outside distinction of the surface. However, in three-dimensional space, twist has an intrinsic chirality that does not require a global surface orientation. A chosen local reference direction along the strand provides a convention for distinguishing the two chiral senses. This intrinsic chirality of twist in three-dimensional space implies that manifoldness and global surface orientation are not required to define or control twisting and linking behavior. This observation motivates representations that encode local cyclic structure directly, without relying on surface-based assumptions.

We therefore generalize these ideas by extending rotation-based representations beyond manifold constraints. By directly encoding cyclic ordering and adjacency in a more general combinatorial structure, hypergraph rotation systems provide a unified scaffold that retains the determinism of rotation systems while supporting more complex connectivity. This shift enables a more directly controllable design space for linked knot structures without the limitations imposed by graph-only or manifold-based approaches.

Table~\ref{table1} provides a comparison with previous approaches. We want to point out that
Block Mesh~\cite{akleman2015blockmeshes} provides a substantially broader theoretical framework based on graphs embedded in 3-manifolds. Its generality
implicitly permits the topological configurations needed to represent
arbitrary knots and links, although such constructions were neither
discussed nor demonstrated in that work. H-GRS does not aim to exceed
Block Mesh in general representational power. Instead, it isolates the
cyclic edge--face incidence and integer-twist information needed for
linked-knot construction, without requiring an underlying 3-manifold
representation, and develops this structure explicitly into the
constructive framework presented here.


\begin{figure*}[htb!]
    \centering    \includegraphics[width=0.16\textwidth]{images/winged/4_2.jpg}    \includegraphics[width=0.16\textwidth]{images/winged/6_2.jpg}   \includegraphics[width=0.16\textwidth]{images/winged/6_3.jpg}    \includegraphics[width=0.16\textwidth]{images/winged/12_2.jpg}    \includegraphics[width=0.16\textwidth]{images/winged/12_3.jpg}    \includegraphics[width=0.16\textwidth]{images/winged/12_4.jpg}\\
    \parbox[t]{0.16\textwidth}{\centering 4 faces, 2 twist, \\2 loops}
    \parbox[t]{0.16\textwidth}{\centering 6 faces, 2 twist, \\2 loops}
    \parbox[t]{0.16\textwidth}{\centering 6 faces, 3 twist, \\3 loops}
    \parbox[t]{0.16\textwidth}{\centering 12 faces, 2 twist, \\2 loops}
    \parbox[t]{0.16\textwidth}{\centering 12 faces, 3 twist, \\3 loops}
    \parbox[t]{0.16\textwidth}{\centering 12 faces, 4 twist, \\4 loops}
    \caption{For local connectivity relations, we can change the number of twists to connect different loops and have disconnected loops at the end. Here we show examples of an edge shared by multiple faces and changing the number of twists to get various connectivity relations.}
    \Description{Six examples show edges shared by different numbers of faces, with varying twist values producing different numbers of linked loops. Strand colors distinguish the resulting components.}
    \label{fig:twistingloops}
\end{figure*}




\section{Topological Representation}
\label{sec:representation}

Our framework is based on the observation that complex linked knot structures can be generated from simple scaffolds using only local combinatorial rules. In particular, we define the combinatorial cycle structure through cyclic orderings of faces around edges together with integer-valued edge labels, allowing local assignments to determine global strand connectivity. This representation is minimal, does not require consistent surface orientation, and applies uniformly to meshes with arbitrary connectivity.

We consider a polygonal mesh as a combinatorial scaffold consisting of faces and edges. At each edge, the incident faces are arranged in a cyclic order, capturing the local radial structure without imposing a global orientation. Each edge is additionally assigned an integer-valued twist, which encodes how adjacent faces are reconnected across the edge. Together, these two ingredients define a local rule that propagates across the mesh to produce globally consistent linked structures.





\subsection{Scaffold Mesh as a Hypergraph Rotation System (H-GRS)}

We distinguish between the input scaffold and the resulting linked geometry. The scaffold encodes only combinatorial structure; the geometric realization of strands is constructed subsequently.
The guarantees provided by the H-GRS representation are combinatorial: the cyclic orderings and twist assignments determine the connectivity and decomposition of the induced strand cycles. They do not, by themselves, determine equivalence between arbitrary spatial knot embeddings or guarantee collision-free geometry. The geometric realizations produced by our construction should instead be viewed as starting points that may require further geometric processing depending on the desired application.

We model the scaffold as a \emph{hypergraph rotation system}, which extends classical rotation systems from graphs to non-manifold, higher-order incidence structures. The essential extension is not higher edge valence alone, but the explicit cyclic ordering of edge--face incidences around each radial edge, together with the cyclic ordering of edges around each face. Accordingly, H-GRS is not intended to represent arbitrary non-manifold topology, but rather the restricted class of edge--face incidence structures required for linked-knot construction.

\begin{definition}[Hypergraph Rotation System / Scaffold Mesh]
\label{def:HGRS}
A \emph{scaffold mesh}, or equivalently a \emph{hypergraph rotation system (H-GRS)}, is an ordered incidence structure
\[
M = (F, E, \{\sigma_{\vec{e}}\}_{\vec{e} \in \vec{E}}, \{\rho_f\}_{f \in F}),
\]
where:
\begin{itemize}
    \item $F$ is a finite set of faces,
    \item $E$ is a finite set of edges, and $\vec{E}$ denotes the set of the two reference orientations of these edges,
    \item for each edge $e \in E$, $F_e \subseteq F$ denotes the set of faces incident to $e$,
    \item for each face $f \in F$, $E_f \subseteq E$ denotes the set of edges bounding $f$,
\end{itemize}
together with:
\begin{itemize}
    \item for each reference orientation $\vec{e} \in \vec{E}$ of an edge, a cyclic ordering $\sigma_{\vec{e}}$ of $F_e$, representing the oriented radial ordering of faces around $e$ relative to the chosen direction $\vec{e}$,
    \item for each face $f \in F$, an \emph{unoriented cyclic ordering} $\rho_f$ of $E_f$, representing the boundary ordering of edges around face $f$, where an ordering and its reversal are identified.
\end{itemize}

The cyclic ordering associated with the opposite direction satisfies
\[
\sigma_{\overleftarrow{e}} = \sigma_{\vec{e}}^{-1}.
\]

No manifold condition or global surface orientation is assumed. In particular, an edge may be incident to an arbitrary number of faces, and a face may have an arbitrary number of edges.

The cardinality $|F_e|$ is called the \emph{valence} of edge $e$, and $|E_f|$ is the \emph{degree} of face $f$.
\end{definition}


\textit{Consistency and orientation of cyclic orderings.}
The H-GRS structure defines its cyclic orderings combinatorially, without requiring manifoldness or a global surface orientation. The face-boundary ordering $\rho_f$ is unoriented: a cyclic ordering and its reversal represent the same face boundary structure. In contrast, the radial cyclic ordering at each edge is oriented, although the scaffold edge itself is undirected. A direction $\vec{e}$ is chosen solely as a local reference to represent this orientation. Reversing the reference direction replaces $\sigma_{\vec{e}}$ by its inverse $\sigma_{\overleftarrow{e}}=\sigma_{\vec{e}}^{-1}$.

%%%%%%%%%%%%%

When the scaffold is realized in three-dimensional space, the prescribed radial cyclic ordering must be geometrically consistent with the arrangement of the incident faces around each edge. A chosen reference direction along the realized edge determines a two-dimensional perpendicular plane in which the two radial senses can be distinguished. The combinatorial ordering $\sigma_{\vec{e}}$ is associated with one of these geometric senses, while reversing the reference direction reverses both the geometric radial sense and the corresponding cyclic ordering. This same convention provides a consistent way to identify the two chiral senses of twist.

Thus, edge direction is auxiliary rather than intrinsic to the scaffold. The oriented radial cyclic ordering is part of the H-GRS data, while its clockwise or counterclockwise interpretation arises from a geometrically consistent three-dimensional realization. Together with the face-boundary orderings, these radial orderings define a generalized rotation system that encodes the local topology of the scaffold without requiring manifoldness or global orientation.

\subsection{Knotted Surface Representation on H-GRS}

Given a scaffold mesh represented as a hypergraph rotation system
$M = (F,E,\{\sigma_{\vec e}\}_{\vec e\in \vec E},\{\rho_f\}_{f\in F})$, we define the connectivity of the knotted surface represented by the H-GRS through integer-valued twist assignments on edges.

\begin{definition}[Knotted Surface Representation]
Let $M$ be a scaffold mesh (H-GRS). A \emph{knotted surface representation}
is a pair $(M, t)$, where
\[
t : E \to \mathbb{Z}
\]
assigns an integer twist to each edge.
\end{definition}

A \emph{face-side} is an incidence of a face with one of its bounding edges; thus, each edge $e$ has one incident face-side for every face in $F_e$. For a chosen local reference direction, we write $\sigma_e$ for the corresponding oriented radial ordering. For each edge $e \in E$ with valence $m = |F_e|$, the cyclic ordering $\sigma_e$ defines an indexing of the incident faces:
\[
\sigma_e = (f_0, f_1, \ldots, f_{m-1}).
\]
The twist $t(e)$ acts as a cyclic shift on this ordering:
\[
f_i \mapsto f_{i + t(e) \; (\mathrm{mod}\; m)}.
\]

This action determines how face-side segments are reconnected across edge $e$. Applying these local reconnections over all edges defines a global adjacency relation on face-sides.

\textit{Induced cycles.}
The resulting adjacency induces a set of cycles obtained by tracing face-side connections across edges according to the twist assignments. Each cycle corresponds to a connected sequence of face-sides obtained by concatenation across edges.

Equivalently, twist assignments can be interpreted as gluing faces 
along edges with a cyclic shift, reordering their boundary adjacencies. 
This operation concatenates face boundaries into longer cycles, which 
define the filaments (strands) of the resulting knotted surface associated with $(M, t)$.
This interpretation is consistent with earlier work, in which twist operations were viewed as concatenating face boundaries via local reconnections~\cite{akleman2009cyclic,akleman2015extended,akleman2020topologically,akleman2015blockmeshes}.



\textit{Trivial representation.}
The special case $(M, 0)$, where $t(e) = 0$ for all $e \in E$, induces no reconnections. In this trivial case $(M,0)$, each face forms an independent cycle, and the resulting surface coincides with the scaffold.  Once such a consistent orientation is fixed on $M$, all nonzero twist assignments $(M,t)$ produce well-defined global cycle structures through
the same local rules.

\textit{Cycle structure under twists.}

The action of twist assignments on a single edge induces a cyclic
group structure that determines the resulting connectivity.

\begin{theorem}[Local cycle structure at a single edge]
\label{theo:twist}
Let $e \in E$ be an edge of valence $m$, with cyclic ordering
$\sigma_e = (f_0, f_1, \ldots, f_{m-1})$. Let $t(e) = k \in \mathbb{Z}$
be the twist assigned to $e$.

Then the twist acts as the cyclic permutation
\[
f_i \mapsto f_{i+k \bmod m}.
\]

The induced cycle structure consists of exactly
\[
\gcd(m, k)
\]
disjoint cycles, each of length
\[
\frac{m}{\gcd(m, k)}.
\]

In particular:
\begin{itemize}
\item If $\gcd(m, k) = 1$, the action produces a single cycle
(unicursal case).
\item If $k \equiv 0 \ (\mathrm{mod}\ m)$, the action is trivial and
each face remains an independent cycle.
\item Otherwise, the faces are partitioned into multiple disjoint
cycles of equal length.
\end{itemize}
The global cycle structure of $(M,t)$ is obtained by composing these local permutations across all edges via the traversal of Algorithm~\ref{alg:corner_graph}.
\end{theorem}



\begin{proof}
Let the $m$ incident face positions be indexed cyclically by $\mathbb{Z}_m$. A twist by $k$ induces the local permutation $i\mapsto i+k \pmod m$. If $k\equiv 0\pmod m$, this permutation is the identity, so every position forms a singleton orbit. Using the standard convention $\gcd(m,0)=m$, this case is consistent with the general formula, giving $m$ orbits of size $1$. Otherwise, the orbit of any position is obtained by repeatedly adding $k$ modulo $m$. These orbits are the cosets of the subgroup generated by $k$ in $\mathbb{Z}_m$, which has $m/\gcd(m,k)$ elements and therefore $\gcd(m,k)$ cosets. Hence the local permutation consists of exactly $\gcd(m,k)$ orbits, each containing $m/\gcd(m,k)$ positions. In particular, if $\gcd(m,k)=1$, all $m$ positions belong to a single local orbit. These are local orbit counts around the edge; the global LK cycles are obtained by composing the local connections throughout the H-GRS.
\end{proof}







This result explains the connectivity patterns observed in Figures~\ref{fig:2_mantwist},~\ref{fig:3_mantwist}, and~\ref{fig:twistingloops}, where varying the twist value produces different cycle decompositions depending on its relation to the edge valence.

\textit{Dynamic topology induced by twists.}
A key feature of the proposed representation is that the hypergraph rotation system does not define a fixed surface. Instead, it provides
a combinatorial scaffold on which topology is induced through local operations.

As explained above, the cyclic orderings $\{\sigma_e\}$ and $\{\rho_f\}$ define only the
\emph{potential} adjacency between faces. The actual connectivity is determined by the twist function $t$, whose action at each edge modifies
the local incidence structure by cyclically reindexing the faces around that edge.

As a result, assigning twist values does not merely annotate the scaffold: it actively transforms it. Faces may be concatenated into longer cycles,
split into multiple components, or rearranged depending on the twist
values. Thus, different assignments $t$ induce different global
topologies over the same underlying H-GRS.

This perspective is central to our framework: topology is not encoded
directly in the scaffold, but emerges from the interaction between the
hypergraph rotation system and the twist assignments.



\subsection{Cycle Structures and Face Operations}

The action of twists on the cyclic orderings induces a combinatorial
process on faces, which can be understood in terms of merging, splitting,
and reordering operations.

\begin{definition}[Cycle Structure]
Let $(M, t)$ be a knotted surface representation. The induced connectivity yields a set of cycles obtained by tracing face-side adjacency across edges according to the twist assignments.
\end{definition}

Equivalently, this process can be interpreted as gluing faces along edges with a cyclic shift, reordering their boundary adjacencies, thereby concatenating boundaries into cycles. These cycles provide a complete description of
the resulting surface topology.

From this viewpoint, twist assignments act as local operators on the
H-GRS that modify cycle structure. In particular:
\begin{itemize}
    \item twists may \emph{merge} multiple face cycles into a single cycle,
    \item twists may \emph{split} a cycle into multiple components,
    \item twists may \emph{permute} the ordering of face adjacency within cycles without changing the
    number of components.
\end{itemize}

These local operations collectively determine the global cycle structure, as described explicitly in Section~3.4.

\begin{figure}[htb!]
    \centering
    \begin{subfigure}[b]{0.46\columnwidth}  
    \includegraphics[width=0.46\textwidth]{images/vertex_hinge/tet_vertex.jpg}    
    \includegraphics[width=0.46\textwidth]{images/vertex_hinge/tet_vertex_spiral.jpg}
        \caption{Vertex hinged Non-Manifold mesh and corresponding LKs.}
        \label{hinge0}
    \end{subfigure}
    \hfill
        \begin{subfigure}[b]{0.46\columnwidth}
        \includegraphics[width=0.46\textwidth]{images/vertex_hinge/tet_edge.jpg}
        \includegraphics[width=0.46\textwidth]{images/vertex_hinge/tet_edge_spiral.jpg}
        \caption{Edge (Piano) hinged mesh and corresponding LKs.}
        \label{hinge1}
    \end{subfigure}
    \hfill
  \caption{The H-GRS structure, by design, does not support hinged vertices. The example on the left shows that non-manifold structures connected only at a vertex do not produce connected LK-structures. The corresponding H-GRS is also not connected, since vertex-only adjacency does not induce connectivity in the H-GRS. In contrast, if the scaffold H-GRS is connected, the resulting structure is guaranteed to be linked by assigning non-zero labels, as in the example on the right.}
  \Description{Two examples compare vertex-only and edge-connected non-manifold scaffolds with their resulting linked-knot structures. Vertex-only adjacency produces separate components, while edge adjacency produces an interlinked structure.}
  \label{fig:hinge_types}
\end{figure}

\textit{Remark on vertex adjacency.}
In the H-GRS framework, vertices are not treated as explicit topological entities, and connectivity is defined solely through edge--face incidences. In particular, adjacency through a shared vertex does not induce connectivity. Consequently, non-manifold configurations that are connected only at a vertex are not connected as H-GRS structures and do not produce connected LK-structures. This distinction, as illustrated in Figure~\ref{fig:hinge_types}, is important. Note that if the scaffold H-GRS is connected, assigning non-zero labels ensures that the resulting strands interact at every edge, producing a linked structure (though not necessarily a single connected component as shown in Figure~\ref{hinge1}). In contrast, the scaffold in Figure~\ref{hinge0} is
not linked as an H-GRS, since vertex-only adjacency does not induce connectivity in the edge--face incidence structure. This is why vertices are not explicitly included in the Definition~\ref{def:HGRS}.

\subsection{Twist-Induced Connectivity and Cycle Extraction}

The central operation in our construction is the reassignment of connections between face-sides incident to each edge.

Consider an edge $e$ that is incident to $m$ face-sides with a cyclic radial ordering. We conceptually replace a small neighborhood of $e$ by a prism-like connector whose side strips correspond to these incident face-sides. A twist label $t(e)$ specifies a \emph{discrete twist} applied along the length of this connector.

This operation should not be interpreted as a rigid geometric rotation. Instead, it is a \emph{helical reconnection}: a face-side entering the connector at one end exits at the face-side shifted by $t(e)$ positions in the cyclic ordering at the other end. Positive and negative values of $t(e)$, relative to the chosen local reference direction, correspond to opposite chiral twists. Thus, each edge acts as a local twisting element that re-routes strands according to its twist label.

\textit{Induced permutation.}
The twist operation at each edge defines a local permutation on its incident face-sides. Combined with the natural adjacency along face boundaries, these local rules induce a global permutation over the set of all face-sides.

\textit{Cycle extraction.}
The strands of the resulting LK structure correspond to the orbits (cycles) of this global permutation.

\textit{Algorithm.}
\begin{enumerate}
    \item Represent each face-side (i.e., each incidence of a face with one of its edges) as a directed element.
    \item Initialize all face-sides as unvisited.
    \item For each unvisited face-side:
    \begin{enumerate}
        \item Start a new strand and mark the current face-side as visited.
        \item Traverse the structure by alternating between:
        \begin{itemize}
            \item \textbf{Face step:} move to the next face-side along the boundary of the current face.
            \item \textbf{Edge step:} apply the twist at the incident edge $e$, sending the current face-side to the one shifted by $t(e)$ positions in the cyclic radial ordering. 
        \end{itemize}
        Alternate these two steps, marking each encountered face-side as visited, until returning to the starting face-side.
        \item Continue until the traversal returns to the starting face-side.
    \end{enumerate}
    \item Each completed traversal defines one strand.
\end{enumerate}

This procedure uniquely determines the number of strands and their connectivity purely from the combinatorial structure of the H-GRS and the twist labels. In particular, for an edge of valence $m$ with twist $t$, the local contribution to strand structure is governed by $\gcd(m, t)$, consistent with Theorem~\ref{theo:twist}.

\section{Construction of Linked Knot Structures}


Given a scaffold mesh represented as a hypergraph rotation system
(H-GRS) and a set of twist assignments, the induced cycle structure
defines the topology of the resulting surface.

We construct linked knot geometries by realizing these cycles as
continuous strands embedded in space. Each cycle obtained from the
knotted surface representation $(M,t)$ corresponds to a closed filament, and the collection of all cycles forms the final linked
structure.

Thus, geometric construction reduces to a realization problem. Given the combinatorial cycles induced by the local twist operators on the H-GRS, we embed them as spatial curves while preserving their connectivity and adjacency relations.

\begin{figure}[htb!]
    \centering        \includegraphics[width=0.23\columnwidth]{images/tetchain/0.png}
\includegraphics[width=0.23\columnwidth]{images/tetchain/1.png}
\includegraphics[width=0.23\columnwidth]{images/tetchain/2.png}
\includegraphics[width=0.23\columnwidth]{images/tetchain/3.png}
  \caption{Four distinct chainmail structures obtained from a regular tetrahedron by assigning twist labels of +2 and -2 to its edges. Each choice produces a different chainmail configuration. In total, even this simplest mesh yields 12 distinct chainmail cases. }
  \Description{Four tetrahedral chainmail structures show distinct interlinking patterns produced by different twist assignments. Colors distinguish the separate loops.}
  \label{fig:tetra_link}
\end{figure}

\subsection{Chainmail Construction}

Chainmail structures arise when twist assignments preserve the
original face decomposition while modifying the local adjacency
between components.

\begin{theorem}[Chainmail Construction]
Let $(M,t)$ be a knotted surface representation. Suppose that for every
edge $e \in E$ of valence $m_e$, the twist label satisfies
\[
t(e) = a_e \, m_e
\]
for some integer $a_e$. Then no edge merges distinct incident faces,
and each original face of the scaffold remains an independent cycle.
\end{theorem}

\begin{proof}
From Section 3, the twist at edge $e$ acts as a cyclic shift on the
ordering $\sigma_e$ of length $m_e$. When $t(e) = a_e m_e$, this shift
is congruent to zero modulo $m_e$, and therefore acts as the identity.

Thus, no local reconnection occurs at any edge, and the induced cycle
structure coincides with the trivial representation $(M,0)$, in which
each face forms an independent cycle.
\end{proof}

Although the cycle structure is unchanged, nonzero twist assignments
modify the geometric realization of adjacent cycles. In particular,
they introduce local interlocking between neighboring components,
producing chainmail configurations.

From the perspective of H-GRS operators, chainmail structures
correspond to twist assignments that preserve the cycle decomposition
while altering the spatial embedding. This separation between
topological invariance and geometric interaction enables the design
of flexible assemblies with independent yet interlocked components.




\subsection{Unicursal (Single Cycle) Knot Construction}
\label{sec:unicursal}
A complementary objective is to merge all face cycles into a single
cycle, producing a unicursal knot or simply a knot.

From Section~3, twists act as local operators on the hypergraph
rotation system, modifying the cycle structure by merging or splitting
face cycles. In particular, a twist at an edge can merge all incident
faces into a single cycle when the twist value is relatively prime to
the local valence.

This observation leads to a simple constructive strategy based on a
spanning tree over faces.

\begin{definition}[Face Adjacency Graph]
Given a scaffold mesh $M=(F,E,\{\sigma_e\},\{\rho_f\})$, the
\emph{face adjacency graph} is the graph whose nodes correspond to
faces in $F$, and where two faces are connected if they are incident
to a common edge.
\end{definition}

\begin{figure}[htb!]
    \centering        \includegraphics[width=0.23\columnwidth]{images/tetknot/0.png}
\includegraphics[width=0.23\columnwidth]{images/tetknot/1.png}
\includegraphics[width=0.23\columnwidth]{images/tetknot/2.png}
\includegraphics[width=0.23\columnwidth]{images/tetknot/3.png}
  \caption{Four distinct unicursal knots obtained from a regular tetrahedron by twisting edges +1, -1, +2, and -2 to achieve a different single cycle in each case. In total, even this simplest symmetric mesh yields 32 distinct unicursal knot cases.}
  \Description{Four single-loop knot structures derived from the same tetrahedral scaffold show distinct winding patterns produced by different twist assignments.}
  \label{fig:tetra_knot}
\end{figure}



Before proving the general case, we first consider the special case in which every edge has valence one or two. This result extends our earlier theorem for 2-manifold structures, where the allowed twist values were restricted to $0$ and $+1$ \cite{xing2010single}. Here, the same construction applies to arbitrary integer twists: any even twist has the same connectivity effect as $0$, while any odd twist has the same connectivity effect as $+1$.

\begin{theorem}[Unicursal Construction for Valence-Two H-GRS] Let $M$ be a connected H-GRS in which every edge has valence one or two. Then there exists an integer twist assignment that produces a unicursal LK. In particular, any odd integer may be assigned to the valence-two edges of a spanning tree of the face-adjacency graph, and any even integer may be assigned to the remaining valence-two edges. \end{theorem}

\begin{proof} Valence-one edges do not affect cycle connectivity and may therefore be ignored. Let $G_F$ denote the face-adjacency graph whose vertices correspond to faces of $M$ and whose edges correspond to valence-two H-GRS edges shared by pairs of faces. Since $M$ is connected, $G_F$ is connected and therefore contains a spanning tree $T$. Root $T$ at an arbitrary face and process its edges outward from the root, so that each processed tree edge connects the already processed subtree to one previously unprocessed face.

Starting from the zero-twist configuration, every face forms a distinct cycle. At a valence-two edge, an even twist preserves the local pairing, whereas an odd twist exchanges the pairing of the two incident strands. Consequently, when an odd twist is applied to a tree edge, the cycle corresponding to the newly added face is merged with the single cycle formed by the already processed subtree. By induction, after processing any connected subtree of $T$, all faces in that subtree belong to one cycle.

Since $T$ spans all faces, after all tree edges have been processed, all faces belong to a single cycle. Assigning even twists to the remaining valence-two edges preserves the corresponding local pairings and therefore does not alter this connectivity. Hence the resulting LK is unicursal. \end{proof}

In the general higher-valence case, a direct constructive procedure based only on the original H-GRS may not always be available. A twist at a radial edge can simultaneously connect and disconnect existing cycles, potentially destroying parts of an already constructed tree. Nevertheless, a unicursal LK can still be obtained by minimally modifying the H-GRS when such obstructions occur. We introduce an \emph{edge-splitting} operation that replaces a problematic radial edge by two geometrically overlapping radial edges and partitions its incident face-sides between them, thereby reducing the valence of each resulting edge. This operation changes the combinatorial H-GRS but preserves its geometric realization. Edge splitting is applied only when necessary and only until the unicursal construction can continue. Since repeated edge splitting can always reduce problematic edges to the valence-two case, for which the preceding theorem provides a constructive solution, this leads to the following general result.

\begin{theorem}[Unicursal Construction] Let $M$ be a connected H-GRS. A unicursal LK can always be constructed from $M$, allowing radial edges to be split when necessary while preserving the geometric realization. \end{theorem}

\begin{proof} Begin with the zero-twist configuration, in which every face forms a distinct cycle, and grow a tree connecting these cycles. At each stage, examine the radial edges on the boundary of the current tree. If a boundary edge admits a twist that connects additional cycles without destroying the tree already constructed, apply such a twist and continue the construction. A higher-valence boundary edge, however, may have no twist that extends the tree without simultaneously reconnecting face-sides already belonging to the same cycle and thereby creating an unwanted cycle. In this case, split the edge by replacing it with two geometrically overlapping radial edges and partitioning its incident face-sides between them. This operation preserves the geometric realization while strictly reducing the valence of the resulting edges.

After each edge split, reconsider the resulting boundary edges and continue with an edge that allows the tree to grow with minimum intervention. If no suitable twist is available, split a problematic boundary edge again. Since every edge split strictly reduces the valence of the affected edge, this process cannot continue indefinitely. In the worst case, all problematic edges can be split until their valence is at most two. The preceding valence-two unicursal construction theorem then guarantees that the remaining face cycles can be connected by a spanning-tree construction. Therefore, the edge-splitting and connection procedure always terminates with a single cycle, and hence produces a unicursal LK. \end{proof}

\noindent \textbf{Remark} An edge split must respect the cyclic order of the face-sides incident to the original radial edge. Specifically, each resulting edge receives a contiguous subsequence of the original cyclic ordering, so that the split preserves the local ordering and geometric realization of the H-GRS. 

\subsection{Design Space and Combinatorial Richness}

The constructions above define a large discrete design space. For a
fixed scaffold with edge set $E$ and a finite set of $k$ allowed twist
values, there are $k^{|E|}$ possible edge-twist assignments before
accounting for symmetry. Since the numerical values of the twists are
irrelevant for this enumeration, the $k$ distinct twist values can be
treated as $k$ distinct edge labels. Assignments related by
orientation-preserving symmetries of the scaffold are considered
equivalent, while mirror-related assignments remain distinct because
reflection reverses twist chirality.

The resulting symmetry-reduced design space can be counted using
Burnside's lemma \cite{burnside1911theory}. If $G$ is the orientation-preserving symmetry group
of the scaffold and $c_E(g)$ is the number of cycles in the permutation
of edges induced by $g\in G$, then the number of symmetry-distinct
assignments is
\begin{equation}
N_G(k)=\frac{1}{|G|}\sum_{g\in G} k^{c_E(g)}.
\end{equation}
A complete derivation and enumeration for the representative
scaffolds used in this work are provided in the Supplementary
Material.

This analysis reveals two distinct growth behaviors. For a fixed
scaffold, $N_G(k)$ is a polynomial in $k$ of degree $|E|$. In contrast,
for a fixed number $k>1$ of available twist values, the underlying
design space grows exponentially with the number of scaffold edges.
Quotienting by scaffold symmetry removes equivalent assignments but
does not eliminate this exponential dependence. For example, with
three distinct twist values, the numbers of symmetry-distinct
configurations are 87 for a tetrahedron, 22,815 for a cube, and
approximately $3.43\times10^{12}$ for a dodecahedron.

These symmetry-distinct assignments represent different topological
configurations of the H-GRS under the chosen equivalence relation.
The enumeration considers all resulting LK configurations together,
without separating them according to the number of components or
whether they belong to particular subclasses such as chainmails or
unicursal knots. Particular subclasses can be obtained by imposing
additional constraints on the twist assignments. For example, in the
two-fold case, restricting the allowed labels to nonzero even twists
produces cycle-preserving chainmail configurations. Unicursal
configurations are more constrained: producing a single global cycle
requires a suitable set of cycle-merging edges, as described by the
spanning-tree construction in Section~\ref{sec:unicursal}.


\begin{figure}[htb!]
    \centering
    \begin{subfigure}[b]{0.32\columnwidth}       \includegraphics[width=\textwidth]{images/eventwisting/cube_1_0.jpg}
        \caption{Only +1 and zero twist}
        \label{singlecycle0}
    \end{subfigure}
    \hfill
    \begin{subfigure}[b]{0.32\columnwidth}
        \includegraphics[width=\textwidth]{images/eventwisting/cube_1_2.jpg}
        \caption{Replacing zero twists with 2-twists.}
        \label{singlecycle1}
    \end{subfigure}
    \hfill
        \begin{subfigure}[b]{0.32\columnwidth}
        \includegraphics[width=\textwidth]{images/eventwisting/cube_1_4.jpg}
        \caption{Replacing zero twists with 4-twists.}
        \label{singlecycle2}
    \end{subfigure}
    \hfill
  \caption{An example of further refinement of unicursal structures. Replacing zero-twist edges with any non-zero even twist in 2-manifold meshes still results in single-cycle constructions. While zero twists produce parallel strands that are only weakly constrained geometrically, non-zero even twists preserve the zero-twist effect at the combinatorial level while introducing additional crossings that locally bind strands together. This demonstrates how multi-twist labels increase geometric coherence without altering topological guarantees.}
  \Description{Three single-cycle structures derived from the same scaffold show increasing local twisting while preserving the same overall connectivity.}
  \label{fig:single_cycle}
\end{figure}

\textit{Further refinement of unicursal structures.}
In the unicursal construction, twist assignments along a spanning tree
are used to merge all face cycles into a single cycle. The remaining
edges, which are not part of the spanning tree, provide additional
degrees of freedom that do not affect the global topology (Figure~\ref{fig:single_cycle}).

In particular, for any edge $e \notin T$ with valence $m_e$, assigning
a twist of the form
\[
t(e) = a_e \, m_e
\]
induces the identity action on the cyclic ordering $\sigma_e$.
Therefore, such twists do not alter the cycle structure and preserve
the unicursal property.

These additional twists can be used to modify the geometric realization
of the structure without changing its connectivity. In practice, they
enable the design of configurations in which the resulting strand
appears more uniformly connected, avoiding loosely attached or
``flapping'' segments.

Even under this constraint, the number of valid assignments remains
large, as each non-tree edge admits multiple independent choices.
Thus, unicursal structures retain a rich design space, combining
topological control through the spanning tree with geometric freedom
through the remaining edges.


\begin{figure}[htb!]
    \centering
    \includegraphics[width=0.24\linewidth]{images/unfolding/box.jpg}
    \includegraphics[width=0.24\linewidth]{images/unfolding/box_tw.jpg}
    \includegraphics[width=0.24\linewidth]{images/unfolding/box_closed.jpg}
    \includegraphics[width=0.24\linewidth]{images/unfolding/box_extreme.jpg}\\
    \parbox[t]{0.24\linewidth}{\centering Input Mesh}
    \parbox[t]{0.24\linewidth}{\centering Open Lid}
    \parbox[t]{0.24\linewidth}{\centering Closed Lid}
    \parbox[t]{0.24\linewidth}{\centering Open Lid}
    \caption{An example of Piano-Hinged systems. Using unfolded polyhedra, we can design joints that can rotate around an edge. In this example, we are using an unfolded face of a cube to design the lid.}
    \Description{Four views show an unfolded cube scaffold and the corresponding linked-knot structure with a hinged lid in different rotational positions.}
    \label{fig:lidexample}
\end{figure}

\section{Algorithmic Realization of the Framework}
\label{sec:algorithmic_realization}

This section presents the algorithmic procedures for constructing linked knot structures from the combinatorial representation described in Section~\ref{sec:representation}. We first extract the cycle structure induced by twist assignments, then construct a connectivity graph based on corner vertices, and finally generate the geometric realization as spatial curves.

\subsection{Corner Connectivity Graph Construction}

Given the cycle structure, we construct an intermediate graph representation based on corner vertices. Each corner represents a face-vertex incidence and is positioned geometrically using local averaging. The twist assignments then determine how corners are connected across edges.

\textit{Explanation of key operations.}
The corner construction creates intermediate vertices by blending each face vertex with its neighbors along the face boundary. The parameter $\beta$ controls how close corners are to the original vertices versus the edge midpoints.

The frame construction establishes a local coordinate system at each edge. The frame selector $s_e$ chooses a reference vertex (typically from an incident face) that defines the zero-angle direction in the radial ordering. This ensures consistent angular measurements around the edge.

The twist application connects corner $c_i$ of face $f_i$ to corner $c_j$ of face $f_j$, where $j = (i + t(e)) \bmod m_e$. This implements the cyclic shift operation defined by the twist function $t$.

\subsection{Geometric Realization with Helical Curves}


The final step constructs spatial curves by combining helical segments around scaffold edges with smooth transitions through the corner regions
of each face. The helical segments realize the twist assignments, while the corner transitions connect consecutive edge segments according to the
cycle structure.

\textit{Helical curve construction.}
Each curve $\gamma_{e,i}(s)$ combines three components:
\begin{itemize}
    \item Linear interpolation $(1-s)v_1 + sv_2$ along the the chosen local reference direction for the edge
    \item Radial offset $r[\cos\theta(s)\mathbf{n}_1 + \sin\theta(s)\mathbf{n}_2]$ in the perpendicular plane
    \item Angular function $\theta(s) = 2\pi(i/m_e + t(e) \cdot s)$ encoding both initial position and twist rate
\end{itemize}

The initial angle $2\pi i/m_e$ distributes the $m_e$ strands uniformly around the edge. The term $2\pi t(e) \cdot s$ introduces $t(e)$ full rotations along the edge length, creating the helical twist geometry that realizes the combinatorial reconnection.

In practice, we realize these connections by generating scaffold-guided
samples along the edge and corner regions and approximating the resulting
sample sequences with parametric curves. Further implementation details,
including the treatment of unoriented cycles and strand thickening, are
provided in the Supplementary Material.

\begin{algorithm}[htb!]
\caption{Twisted Corner Connectivity Graph Construction}
\label{alg:corner_graph}
\begin{algorithmic}[1]
\Require Mesh $M = (F, E, \{\sigma_{\vec{e}}\}, \{\rho_f\})$ with vertex positions $V$, twist function $t : E \to \mathbb{Z}$, frame selectors $\{s_e\}_{e \in E}$
\Ensure Corner connectivity graph $G_c = (C, A)$

\State $C \gets \emptyset$, $A \gets \emptyset$
\State $\mathcal{M} \gets$ empty mapping from $(f, v)$ pairs to corners

\Function{BuildCorners}{$F, V$}
    \For{each face $f \in F$ with boundary ordering $\rho_f = (v_0, v_1, \dots, v_{n-1})$}
        \For{$i = 0$ to $n-1$}
            \State $v \gets v_i$
            \State $v^- \gets v_{(i-1 \bmod n)}$, $v^+ \gets v_{(i+1 \bmod n)}$
            \State $p \gets \beta \cdot v + \frac{1-\beta}{2}(v^- + v^+)$ \Comment{Corner position}
            \State Create corner vertex $c_{f,i}$ at position $p$
            \State $C \gets C \cup \{c_{f,i}\}$
            \State $\mathcal{M}(f, v_i) \gets c_{f,i}$
        \EndFor
    \EndFor
\EndFunction

\Function{BuildFrame}{$e = (v_1, v_2), s_e$}
    \State $\mathbf{d} \gets \frac{v_2 - v_1}{\|v_2 - v_1\|}$ \Comment{Chosen local reference direction for the edge}
    \State $v_r \gets$ vertex chosen by frame selector $s_e$
    \State $\mathbf{r} \gets v_r - v_1$
    \State $\mathbf{n}_1 \gets \mathbf{r} - (\mathbf{r} \cdot \mathbf{d})\mathbf{d}$ \Comment{Project to orthogonal plane}
    \State Normalize $\mathbf{n}_1$
    \State $\mathbf{n}_2 \gets \mathbf{d} \times \mathbf{n}_1$ \Comment{Complete orthonormal basis}
    \State \Return $(\mathbf{d}, \mathbf{n}_1, \mathbf{n}_2)$
\EndFunction

\Function{ProcessEdge}{$e = (v_1, v_2)$}
    \State $(\mathbf{d}, \mathbf{n}_1, \mathbf{n}_2) \gets$ \Call{BuildFrame}{$e, s_e$}
    \State $F_e \gets$ set of faces incident to edge $e$
    
    \For{each $f \in F_e$}
        \State Compute face centroid $c_f$
        \State Project $(c_f - v_1)$ to plane orthogonal to $\mathbf{d}$
        \State Compute angular coordinate $\theta_f$ in $(\mathbf{n}_1, \mathbf{n}_2)$ plane
    \EndFor
    
    \State Index $F_e$ according to the prescribed radial ordering
$\sigma_e=(f_0,f_1,\ldots,f_{m_e-1})$, with increasing $\theta_f$
consistent with this ordering.
    
    \For{$i = 0$ to $m_e - 1$}
        \State $j \gets (i + t(e)) \bmod m_e$ \Comment{Apply twist}
        \State $c_i \gets \mathcal{M}(f_i, v_1)$ \Comment{Corner of $f_i$ at $v_1$}
        \State $c_j \gets \mathcal{M}(f_j, v_2)$ \Comment{Corner of $f_j$ at $v_2$}
        \State $A \gets A \cup \{(c_i, c_j)\}$ \Comment{Connect twisted corners}
    \EndFor
\EndFunction

\State \Call{BuildCorners}{$F, V$}
\For{each edge $e \in E$}
    \State \Call{ProcessEdge}{$e$}
\EndFor

\State \Return $G_c = (C, A)$
\end{algorithmic}
\end{algorithm}

\begin{algorithm}[htb!]
\caption{Geometric Realization with Combinatorial Frames}
\label{alg:geometric_realization}
\begin{algorithmic}[1]
\Require Mesh $M = (F, E, \{\sigma_{\vec{e}}\}, \{\rho_f\})$ with vertices $V$, corner graph $G_c = (C, A)$, twist function $t : E \to \mathbb{Z}$, frame selectors $\{s_e\}_{e \in E}$, radius parameter $r > 0$
\Ensure Set of spatial curves $\Gamma$

\State $\Gamma \gets \emptyset$

\For{each edge $e = (v_1, v_2) \in E$ }
    \State $(\mathbf{d}, \mathbf{n}_1, \mathbf{n}_2) \gets$ \Call{BuildFrame}{$e, s_e$}
    \State $(f_0,f_1,\ldots,f_{m_e-1}) \leftarrow$ faces around $e$ in the prescribed radial ordering $\sigma_e$
    
    \For{$i = 0$ to $m_e - 1$}
        \State Define parametric curve $\gamma_{e,i} : [0,1] \to \mathbb{R}^3$:
        \Statex \hspace{2em}
        $\displaystyle
        \gamma_{e,i}(s) = (1-s)v_1 + s v_2 + r \left[\cos\theta(s) \, \mathbf{n}_1 + \sin\theta(s) \, \mathbf{n}_2\right]
        $
        \Statex \hspace{2em}
        where $\theta(s) = 2\pi\left(\frac{i}{m_e} + t(e) \cdot s\right)$
        
        \State $\Gamma \gets \Gamma \cup \{\gamma_{e,i}\}$
    \EndFor
\EndFor

\For{each edge $(c_a, c_b) \in A$}
    \State Connect corresponding edge-curve samples through the corner region according to the extracted cycle.
\EndFor

\State \Return $\Gamma$
\end{algorithmic}
\end{algorithm}



\section{Behavior Design}

The proposed framework not only defines topology but also provides a systematic basis for exploring mechanical behavior. In particular, the cycle structures induced by different twist assignments can exhibit qualitatively different behaviors, including compliant and articulated motion (see Figure~\ref{fig:cube_behaviour}). This observation suggests that cycle structure can be used as a design parameter to separate compliance (within regions) from articulation (between regions). We exploit the observation shown in Figure~\ref{fig:cube_behaviour} to construct hybrid structures by decomposing the scaffold into regions with controlled cycle properties and connecting them through hinge-like edges.

\begin{figure}[htb!]
    \centering
    \begin{subfigure}[b]{0.32\columnwidth}       \includegraphics[width=\textwidth]{images/photos/cube_1.jpg}
        \caption{+1 twisted edges.}
        \label{cube1_beh}
    \end{subfigure}
    \hfill
    \begin{subfigure}[b]{0.32\columnwidth}
        \includegraphics[width=\textwidth]{images/photos/cube_2.jpg}
        \caption{+2 twisted edges.}
        \label{cube2_beh}
    \end{subfigure}
    \hfill
        \begin{subfigure}[b]{0.32\columnwidth}
        \includegraphics[width=\textwidth]{images/photos/cube_3.jpg}
        \caption{+3 twisted edges.}
        \label{cube3_beh}
    \end{subfigure}
    \hfill
  \caption{
  \textbf{Comparison of the behavior of three 3D-printed LK structures generated from a cubic scaffold with different edge-twist assignments.} The first structure in \ref{cube1_beh} (twist value 1) is freely moving, while the second in \ref{cube2_beh} is a chainmail configuration that exhibits compliant (spring-like) behavior in our fabricated example. The third structure in \ref{cube3_beh} also exhibits compliant behavior. Notably, the first and third structures have the same number of cycles but differ in their linkage patterns and exhibit distinct mechanical responses. These examples illustrate that different twist assignments can result in qualitatively different behaviors of the structure.
  }
  \Description{Three photographs show cubic linked-knot structures generated with twist values 1, 2, and 3.}
  \label{fig:cube_behaviour}
\end{figure}

\subsection{Region-Based Decomposition}

We decompose the scaffold mesh (H-GRS) into subregions such that each subregion is internally unicursal. That is, within each region, twist assignments merge all faces into a single cycle, producing a component that can exhibit compliant behavior.

Adjacent regions are connected through edges that are assigned twists
of the form
\[
t(e) = a_e \, m_e,
\]
where $m_e$ is the valence of edge $e$. As shown in Section~4.1, such
twists induce the identity action on the cyclic ordering $\sigma_e$,
and therefore preserve the cycle structure.

In our fabricated examples,these edges exhibit hinge-like behavior that allow relative motion between adjacent regions while approximately maintaining their internal coherence.


Thus, the structure is composed of:
\begin{itemize}
    \item \emph{unicursal regions}, which behave as compliant units, and
    \item \emph{hinge-like} edges, which connect regions without merging
    their cycles.
\end{itemize}

\subsection{Region Connectivity Graph}

To analyze the global behavior, we define a higher-level abstraction.

\begin{definition}[Region Connectivity Graph]
Given a decomposition of the scaffold into unicursal regions, the
\emph{region connectivity graph} is defined as follows:
\begin{itemize}
    \item each region is represented as a vertex,
    \item an edge is placed between two vertices if the corresponding
    regions are connected by a hinge edge.
\end{itemize}
\end{definition}

This graph provides a combinatorial description of the connections between regions and can be used to characterize the organization of locally compliant and globally articulated behavior observed in our examples.

\subsection{Classes of Behavior}

Different classes of mechanical behavior arise depending on the
topology of the region connectivity graph.

\textit{Tree structures (Piano hinge systems).}
If the region-connectivity graph forms a tree, the resulting structure can exhibit open articulated behavior. Each hinge-like connection can provide relative motion between adjacent regions, while the absence of cycles provides a simple organization for such articulated behavior.
This generalizes classical hinge chains and Piano-like articulated LK structures (Figures~\ref{teaser2},~\ref{teaser3},~\ref{teaser4}, and~\ref{fig:lidexample}).

\begin{figure}[htb!]
    \centering    \includegraphics[width=0.46\columnwidth]{images/herringbone/herringbone.jpg}    \includegraphics[width=0.46\columnwidth]{images/herringbone/herringbone_links.jpg}
    \caption{Herringbone origami as an LK structure. Left: A planar polygonal mesh representing a herringbone origami pattern, modeled as a 2-manifold surface with boundary. Right: The corresponding LK structure obtained by assigning twist labels to the mesh edges, transforming the surface pattern into a linked-knot configuration.}
    \Description{A planar herringbone mesh is shown beside its corresponding linked-knot structure, which follows the same repeating pattern as interlinked tubular strands.}
    \label{fig:herringbone}
\end{figure}

\textit{Surface-like structures (origami generalization).}
If the region connectivity graph forms a 2-manifold mesh, the structure resembles a network of hinge-like connections defined over a surface. In this case, regions can behave as panels connected by hinge-like edges, generalizing origami-like systems to non-manifold and higher-valence settings (Figure~\ref{fig:herringbone}).

\textit{Cyclic structures (flexagon-like systems).}
If the region connectivity graph contains closed loops, the structure can exhibit cyclic articulated behavior. Such configurations are analogous to flexagons, where loops of hinge-like connections can support folding and reconfiguration behavior.


\begin{figure}[htb!]
    \centering
    \begin{overpic}[width=0.32\columnwidth]{images/flexagon/isosleces_tet_1.jpg}
        \put(-5,85){1}
    \end{overpic}
    \begin{overpic}[width=0.32\columnwidth]{images/flexagon/isosleces_tet_2.jpg}
        \put(-5,85){2}
    \end{overpic}
    \begin{overpic}[width=0.32\columnwidth]{images/flexagon/isosleces_tet_3.jpg}
        \put(-5,85){3}
    \end{overpic}    \includegraphics[width=0.32\columnwidth]{images/flexagon/isosleces_tet_1_thread.jpg}    \includegraphics[width=0.32\columnwidth]{images/flexagon/isosleces_tet_2_thread.jpg}
   \includegraphics[width=0.32\columnwidth]{images/flexagon/isosleces_tet_3_thread.jpg}
    \begin{overpic}[width=0.32\columnwidth]{images/flexagon/isosleces_tet_4.jpg}
        \put(-5,85){4}
    \end{overpic}
    \begin{overpic}[width=0.32\columnwidth]{images/flexagon/isosleces_tet_5.jpg}
        \put(-5,85){5}
    \end{overpic}
    \begin{overpic}[width=0.32\columnwidth]{images/flexagon/isosleces_tet_6.jpg}
        \put(-5,85){6}
    \end{overpic}
   \includegraphics[width=0.32\columnwidth]{images/flexagon/isosleces_tet_4_thread.jpg}
 \includegraphics[width=0.32\columnwidth]{images/flexagon/isosleces_tet_5_thread.jpg}    \includegraphics[width=0.32\columnwidth]{images/flexagon/isosleces_tet_6_thread.jpg}
    \caption{
    \textbf{An example of cyclic hinges: Six reconfiguration states of a flexagon-based LK structure.} Each image shows a distinct configuration produced by continuous deformation of the same labeled non-manifold surface scaffold, with topology preserved. The structure is defined by integer twist labels on edges, which merge tetrahedral components into knots while exhibiting relative motion between them. This example illustrates how flexagon-like reconfiguration can be realized within the LK framework.
    }
    \Description{Six configurations show continuous reorientation of the same flexagon-based linked-knot structure while preserving its overall topology.}
    \label{fig:flexagon}
\end{figure}


\subsection{Design Implications}

This decomposition provides a unified framework for exploring behavior design:
\begin{itemize}
\item Local twist assignments define the cycle structure of each region and can lead to qualitatively different compliant or articulated behaviors.
\item Hinge-like edges can provide interfaces that allow relative motion between regions.
\item The topology of the region connectivity graph provides a combinatorial basis for organizing different global motion behaviors.
\end{itemize}

By combining these elements, designers can construct structures that integrate different forms of compliance and articulation within a single unified model. Our examples exhibit a range of qualitative behaviors, from flexible continuous strands to articulated mechanisms and hybrid configurations.

\begin{figure}[htb!]
    \centering      
    \includegraphics[width=0.32\columnwidth]{images/photos/flex_1.jpg}
    \hfill
    \includegraphics[width=0.32\columnwidth]{images/photos/flex_2.jpg}
    \hfill
    \includegraphics[width=0.32\columnwidth]{images/photos/flex_3.jpg}
    \hfill
  \caption{Photographs of three different configurations of the 3D-printed articulated LK structure obtained from a flexagon scaffold (see Figure~\ref{fig:flexagon}). The structure moves freely, and the images illustrate its range of motion.}
  \Description{Three photographs show different configurations of the same articulated flexagon-based linked-knot structure.}
  \label{fig:flexon_behaviour}
\end{figure}

\begin{figure*}[htp!]
    \centering
        \begin{overpic}[width=0.23\linewidth]{images/bravais/cp_1.jpg}
            \put(-5,85){\centering 1}
        \end{overpic}
        \begin{overpic}[width=0.23\linewidth]{images/bravais/cp_2.jpg}
            \put(-5,85){\centering 2}
        \end{overpic}
        \begin{overpic}[width=0.23\linewidth]{images/bravais/cp_3.jpg}
            \put(-5,85){\centering 3}
        \end{overpic}
        \begin{overpic}[width=0.23\linewidth]{images/bravais/cp_4.jpg}
            \put(-5,85){\centering 4}
        \end{overpic}
        \includegraphics[width=0.23\linewidth]{images/bravais/cp_1_repeated.jpg}        \includegraphics[width=0.23\linewidth]{images/bravais/cp_2_repeated.jpg}        \includegraphics[width=0.23\linewidth]{images/bravais/cp_3_repeated.jpg}
        \includegraphics[width=0.23\linewidth]{images/bravais/cp_4_repeated.jpg}
    \caption{Using a tileable twisting module, we can use the distinct edges of the cube as a tile. Adjacency of the cubic tessellation gives 4 threads per edge, and there are 3 distinct edges leading to 12 threads per module. Assuming all thread components have the same twisting operations, depending on the twisting number, the connected components change. In the case of 0 and 4 twists (modulo 4 twist), there are 3 distinct components; 1 and 3 twists give 4 distinct components, and 2 twists give 6 distinct components for cubic tessellation.}
    \Description{Four twisting modules and their tiled repetitions show how different uniform twist values change the connectivity of a periodic cubic linked-knot structure.}
    \label{fig:bravais3D_repeated}
\end{figure*}

\section{Triply Periodic Woven and Knitted Structures}

While the previous sections focus on finite constructions and their
articulated or compliant behavior, the proposed framework also extends
naturally to triply periodic settings, enabling the design of volumetric
woven and knitted structures.

\subsection{Periodic H-GRS from 3D Honeycombs}

We consider scaffold meshes derived from triply periodic 3D honeycombs. In this setting, the hypergraph rotation system (H-GRS) is defined over
a periodic domain, and its combinatorial structure repeats under lattice translations. For a periodic H-GRS, twist labels are specified within a fundamental domain and repeated under lattice translations; thus, each labeled edge represents the same twist assignment on all of its translationally equivalent copies. Rather than producing a finite set of cycles, twist assignments induce cycle structures that extend throughout space. These cycles correspond
to strands that propagate across the periodic domain, forming continuous filament networks.

\subsection{Uniform Twists}

We call a periodic twist assignment \emph{uniform} when every edge in the fundamental domain is assigned the same twist value; otherwise, the assignment is \emph{non-uniform}. The same uniform-twist principle extends naturally to 2D and 3D periodic structures; global connectivity is governed entirely by a single, uniformly applied edge-twist parameter.

In Figure~\ref{fig:bravais3D_repeated}, all edges of the cubic tessellation are assigned the same twist label using a tileable twisting module. Because the cube has three distinct edge classes and each edge contributes four thread segments, each periodic module contains twelve threads. As the uniform twist value varies, the number of connected components changes systematically, demonstrating that global connectivity in the cubic tessellation is controlled solely by this single parameter. The same method can be applied to any doubly periodic 2-honeycombs or triply periodic 3-honeycombs, leading to exponentially many planar or volumetric plain woven fabrics.

\begin{figure}
    \centering
    \includegraphics[width=0.32\linewidth]{images/photos/1twist_clear.jpeg}
    \includegraphics[width=0.32\linewidth]{images/photos/2twist_cleaner.jpeg}
    \includegraphics[width=0.32\linewidth]{images/photos/3twist_cleaner.jpeg}
    \caption{Physical realizations of triply periodic linked knot structures derived from the truncated octahedral tessellation. The constructions use a paper scaffold composed of triangular prisms, reflecting the valence-3 edge structure of the tessellation. Left: single-twist configuration. Middle: two-twist configuration. Right: three-twist configuration.}
    \Description{Three photographs show paper realizations of periodic linked-knot structures with one, two, and three twists.}
    \label{fig:troct_paper}
\end{figure}

\subsection{Non-Uniform Twists}

Figures~\ref{fig:bravais2D_repeated_nonuniform}, ~\ref{fig:cp_single}, and ~\ref{fig:cp_multiple} illustrate doubly and triply periodic LK structures generated by assigning \emph{twist labels to distinct edge classes} within a single repeating cell. In contrast to the uniform-twist cases, not all edges of the repeating cell are twisted in the same manner, even though the underlying periodic scaffold remains unchanged. This heterogeneous edge labeling immediately and significantly expands the design space: varying twist values of the edges of cells locally within one fundamental domain is sufficient to produce qualitatively new global interlacing patterns and connectivity types. As these examples demonstrate, once uniformity is relaxed, the space of achievable periodic LK structures grows even more rapidly, enabling the exploration of a large family of previously unexplored configurations even for a single periodic cell.


\begin{figure}[htb!]
    \centering
    \centering
    \begin{overpic}[width=0.24\linewidth]{images/ws_2d/2_1.jpg}
        \put(-5,50){\rotatebox[origin=c]{90}{\centering [2,1]}}
    \end{overpic}
    \includegraphics[width=0.24\linewidth]{images/ws_2d/2_1_tiled.jpg}
    \begin{overpic}[width=0.24\linewidth]{images/ws_2d/2_0_1.jpg}
        \put(-5,50){\rotatebox[origin=c]{90}{\centering [2,0,1]}}
    \end{overpic}
    \includegraphics[width=0.24\linewidth]{images/ws_2d/2_0_1_tiled.jpg}
    \begin{overpic}[width=0.24\linewidth]{images/ws_2d/2_2_0.jpg}
        \put(-5,50){\rotatebox[origin=c]{90}{\centering [2,2,0]}}
    \end{overpic}    \includegraphics[width=0.24\linewidth]{images/ws_2d/2_2_0_tiled.jpg}
    \begin{overpic}[width=0.24\linewidth]{images/ws_2d/2_2_1.jpg}
        \put(-5,50){\rotatebox[origin=c]{90}{\centering [2,2,1]}}
    \end{overpic}    \includegraphics[width=0.24\linewidth]{images/ws_2d/2_2_1_tiled.jpg}
    \caption{With non-uniform twisting for distinct edges, we are able to produce different classes of structures such as knits and chainmails.}
    \Description{Four local twist assignments are shown beside their tiled planar repetitions, producing different periodic linked patterns.}  \label{fig:bravais2D_repeated_nonuniform}
\end{figure}

\begin{figure}[htb!]
    \centering
    \begin{overpic}[width=0.24\linewidth]{images/cp_nonuniform/2_1_1.jpg}
        \put(-5,50){\rotatebox[origin=c]{90}{\centering [2,1,1]}}
    \end{overpic}
    \includegraphics[width=0.24\linewidth]{images/cp_nonuniform/2_1_1_tiled.jpg}
    \includegraphics[width=0.24\linewidth]{images/cp_nonuniform/2_1_1_thread.jpg}
    \includegraphics[width=0.24\linewidth]{images/cp_nonuniform/2_1_1_thread_tiled.jpg}\\
    \begin{overpic}[width=0.24\linewidth]{images/cp_nonuniform/n1_0_1.jpg}
        \put(-5,50){\rotatebox[origin=c]{90}{\centering [-1,0,1]}}
    \end{overpic}
    \includegraphics[width=0.24\linewidth]{images/cp_nonuniform/n1_0_1_tiled.jpg}
    \includegraphics[width=0.24\linewidth]{images/cp_nonuniform/n1_0_1_thread.jpg}
    \includegraphics[width=0.24\linewidth]{images/cp_nonuniform/n1_0_1_thread_tiled.jpg}\\
    \begin{overpic}[width=0.24\linewidth]{images/cp_nonuniform/n5_0_1.jpg}
        \put(-5,50){\rotatebox[origin=c]{90}{\centering [-5,0,1]}}
    \end{overpic}
    \includegraphics[width=0.24\linewidth]{images/cp_nonuniform/n5_0_1_tiled.jpg}
    \includegraphics[width=0.24\linewidth]{images/cp_nonuniform/n5_0_1_thread.jpg}
    \includegraphics[width=0.24\linewidth]{images/cp_nonuniform/n5_0_1_thread_tiled.jpg}\\
    \parbox[t]{0.24\linewidth}{\centering Twisted unit}
    \parbox[t]{0.24\linewidth}{\centering Tiled twisted unit}
    \parbox[t]{0.24\linewidth}{\centering Unit thread}
    \parbox[t]{0.24\linewidth}{\centering Tiled threads}
    \caption{Using a combination of different twist numbers, it is possible to get a single unique component structures, top row shows an example where the threads are infinite and repeat in a $4\times2\times2$ domain. The middle row shows an example where the periodicity is contained within a single cycle curve, spanning $2\times1\times2$ domain. This shape creates planar chainmails that are stacked in one direction that are disconnected. Bottom row shows an example where one of the edges in the middle row is modulo twisted (-1 to -5) to form a link connection between different layers of chainmail, leading to a singly linked structure.}
    \Description{Three periodic twist assignments are shown with their tiled structures and representative strand paths, illustrating different global connectivity patterns.}
    \label{fig:cp_single}
\end{figure}

\begin{figure}[htb!]
    \centering
    \begin{overpic}[width=0.24\linewidth]{images/cp_nonuniform/1_3_1.jpg}
        \put(-5,50){\rotatebox[origin=c]{90}{\centering [1,3,1]}}
    \end{overpic}
    \includegraphics[width=0.24\linewidth]{images/cp_nonuniform/1_3_1_tiled.jpg}
    \includegraphics[width=0.24\linewidth]{images/cp_nonuniform/1_3_1_thread.jpg}
    \includegraphics[width=0.24\linewidth]{images/cp_nonuniform/1_3_1_thread_tiled.jpg}\\
    \begin{overpic}[width=0.24\linewidth]{images/cp_nonuniform/2_3_2.jpg}
        \put(-5,50){\rotatebox[origin=c]{90}{\centering [2,3,2]}}
    \end{overpic}
    \includegraphics[width=0.24\linewidth]{images/cp_nonuniform/2_3_2_tiled.jpg}
    \includegraphics[width=0.24\linewidth]{images/cp_nonuniform/2_3_2_thread.jpg}
    \includegraphics[width=0.24\linewidth]{images/cp_nonuniform/2_3_2_thread_tiled.jpg}\\
    \parbox[t]{0.24\linewidth}{\centering Twisted unit}
    \parbox[t]{0.24\linewidth}{\centering Tiled twisted unit}
    \parbox[t]{0.24\linewidth}{\centering Unit thread}
    \parbox[t]{0.24\linewidth}{\centering Tiled thread}
    \caption{Different twist numbers enable control of the number of unique components in the repeating structure, the top row shows an example where there are two distinct elements with $1\times2\times1$ repeating domain. Bottom row shows an example with three distinct elements with $1\times4\times1$ repeating domain.}
    \Description{Two periodic twist assignments are shown with their tiled structures and representative strand paths, producing two- and three-component repeating configurations.}
    \label{fig:cp_multiple}
\end{figure}

\subsection{Design Implications}

The introduction of translational periodicity significantly expands the design space of H-GRS–based structures. Instead of restricting the construction to a single-cell unit, we assume a \emph{fundamental domain consisting of multiple cells} that tiles space under a translational symmetry group. Within this domain, each edge of every cell can be assigned an independent twist label, subject only to periodic consistency constraints. An example is provided in Figure~\ref{fig:voronoi_example}.

\textit{Full combinatorial control within periodic domains.}
By allowing independent twist assignments across all edges in a multi-cell fundamental domain, the resulting structure is no longer constrained by uniformity or single-cell symmetry. This provides a rich combinatorial design space in which local modifications propagate globally through periodic tiling, enabling highly diverse classes of woven and knitted volumetric structures.

\textit{Multiple H-GRS structures from a single periodic graph.}
A key implication of this formulation is that a single periodic scaffold graph can generate many distinct H-GRS structures. The same underlying periodic graph admits multiple cyclic orderings and edge-label assignments, each defining a different hypergraph rotation system. Consequently, even when the combinatorial adjacency of the scaffold is fixed, the induced LK structures can vary dramatically in their cycle structure, linking behavior, and geometric realization.

This observation is particularly evident in periodic graphs derived from crystallographic constructions, such as Laves graphs, where different choices of cyclic ordering and twist assignments produce a wide variety of woven and interlaced configurations.
Thus, the mapping from periodic graphs to LK structures is inherently many-to-many, greatly amplifying the expressive power of the framework. An example is provided in Figure~\ref{fig:laves_intro}.

\textit{Unified representation of weaving and knitting in 3D.}
The distinction between weaving and knitting emerges naturally from the connectivity induced by twist assignments:
\begin{itemize}
    \item structures with multiple disjoint cycles correspond to interlaced strand systems (weaving and chainmail),
    \item structures with a single connected cycle correspond to unicursal filament systems (knitting),
    \item intermediate configurations produce hybrid behaviors combining both characteristics.
\end{itemize}
Thus, both classical textiles and more complex architectured materials can be described within the same periodic framework.

\textit{Dependence on fundamental domain size and topology.}
Unlike single-cell constructions, the topology of the resulting structure depends not only on the twist values but also on the size and arrangement of the fundamental domain. Increasing the number of cells in the domain allows the realization of connectivity patterns that are impossible in smaller domains, including long-range linking, layered interlocking systems, and spatially varying weave patterns.

\textit{Emergence of higher-order linking and entanglement.}
Non-uniform twist assignments across the domain can introduce linking between strands that span multiple periodic cells. This enables the construction of complex chainmail-like assemblies, volumetric weaves, and interlocked networks whose global connectivity cannot be inferred from any single local configuration.

\textit{Open questions: unicursal vs.\ chainmail periodic structures.}
A fundamental aspect of periodic designs is the classification of resulting strand systems into unicursal (single-component) or multi-component (chainmail-like) structures. In multi-cell domains, this classification depends on subtle interactions between twist labels, cyclic orderings, and translational symmetry.

Understanding these relationships requires analyzing the induced permutation structure over the periodic domain, leading to connections with group theory and covering space arguments. A complete treatment of these questions—including classification, minimal domains, and controllability of component counts—is beyond the scope of this paper and remains an important direction for future work.




\section{Discussion, Limitations, and Future Work}
\label{sec:discussion}

The proposed framework establishes a unified correspondence between topological scaffold structures and the resulting articulated or filament-based configurations. By encoding connectivity through hypergraph rotation systems and local twist assignments, complex structures---including knots, links, chainmail systems, and triply periodic woven or knitted materials---can be generated within a single combinatorial paradigm.

A key contribution of this work is the separation between \emph{scaffold topology} and \emph{geometric realization}. The same underlying combinatorial structure can give rise to qualitatively different behaviors depending on whether it is interpreted as a finite articulated system or as a periodic filament network. This dual interpretation enables a unified treatment of mechanisms (hinges and compliant systems) and materials (weaving and knitting). In the periodic setting, the framework provides a direct link between cycle structure and classes of volumetric textile-like materials: unicursal constructions correspond to continuous filament systems analogous to knitting, while multi-cycle structures correspond to interlaced strand systems analogous to weaving. Intermediate configurations offer a rich design space for hybrid architectured materials.

\begin{figure}[htb!]
    \centering
    \includegraphics[width=0.30\linewidth]{images/voronoi/sites.jpg}\hfill
    \includegraphics[width=0.30\linewidth]{images/voronoi/sites_repeated.jpg}\hfill
    \includegraphics[width=0.30\linewidth]{images/voronoi/sites_repeated_center_voronoi.jpg}\\
    \includegraphics[width=0.30\linewidth]{images/voronoi/voronoi_result.jpg}\hfill
    \includegraphics[width=0.30\linewidth]{images/voronoi/unique_edges_3_val.jpg}\hfill
    \includegraphics[width=0.30\linewidth]{images/voronoi/unique_edges_4_val.jpg}\\
    \includegraphics[width=0.30\linewidth]{images/voronoi/unique_edges.jpg}\hfill
    \includegraphics[width=0.30\linewidth]{images/voronoi/unique_twists.jpg}\hfill
    \includegraphics[width=0.30\linewidth]{images/voronoi/repeated.jpg}
    \caption{An arbitrary fundamental domain consisting of multiple cells can be constructed using a repeating Voronoi domain. For instance, by placing points at the face centers of a cube, we obtain a cell structure that differs from the standard Wigner--Seitz cell; many other constructions are also possible. The top row illustrates the construction of these cells. The middle row shows the resulting periodic structure and highlights edges with distinct valencies. These distinct edges define a minimal repeating unit. By assigning independent labels to each edge, we can control the resulting connectivity and generate the twisting elements shown in the bottom row.}
    \Description{A sequence of panels shows construction of a periodic Voronoi cell, identification of distinct edge classes, assignment of twist values, and the resulting repeated linked-knot structure.}
    \label{fig:voronoi_example}
\end{figure}

More broadly, this approach emphasizes \emph{design through topology}: local combinatorial rules determine global connectivity, enabling systematic exploration of complex structures without explicitly constructing their geometry. We close by discussing the principal limitations of this framework alongside the future directions they suggest.


\subsection{Limitations and Future Directions}

\textit{Cycle-based characterization.}
The framework relies fundamentally on cycle structures induced by the hypergraph rotation system. As observed in the study of spatial graph embeddings, however, not all forms of entanglement are captured purely through cycles: there exist configurations that are globally tangled without containing nontrivial knotted or linked cycles, and cycle-based descriptions do not form complete invariants of spatial entanglement. Extending the framework to incorporate invariants capable of detecting non-cycle-based tangling---particularly for complex or highly constrained scaffolds---remains an important direction for future work.

\textit{Geometric embedding and optimality.} 
The combinatorial construction specifies connectivity but does not uniquely determine a geometric embedding, and the framework does not currently provide a notion of optimality among the embeddings that realize a given (M,t). Multiple geometrically distinct realizations may correspond to the same combinatorial structure, with different physical or mechanical properties, and there is no intrinsic criterion for selecting minimally tangled, energetically favorable, or fabrication-ready configurations. As in classical knot theory, where quantities such as unknotting number are difficult to compute exactly, defining and computing optimal realizations is open. Future work may focus on optimization methods that produce embeddings satisfying criteria such as minimal tangling, energy minimization, or fabrication feasibility, and on quantitative complexity measures analogous to unknotting and untangling numbers, accessible at least through bounds or approximations. Such mappings from topological designs to specific geometric or physical realizations are orthogonal to the core design contribution presented here.

A direction for future work is to remove the dependence on the three-dimensional geometric interpretation of radial ordering. The present construction uses the two-dimensional plane normal to an edge in 3D to establish the angular ordering required for geometric realization. A braid-based combinatorial formulation may provide an alternative way to encode this information without relying on such a normal plane, potentially allowing analogous constructions in higher-dimensional ambient spaces. 

\textit{Computational exploration of the design space.}
The global behavior of the system emerges from local twist assignments, but exploring the full space of assignments is combinatorially expensive. For large scaffolds or periodic domains, systematic enumeration and classification of all possible structures is computationally infeasible without additional constraints or heuristics. Developing efficient algorithms for sampling, classification, and visualization of large families of scaffold-based structures will be essential for practical design applications.

\textit{Periodic-to-finite transitions.}
The extension to triply periodic structures assumes idealized infinite repetition. In practice, physical realizations are finite and may introduce boundary effects that are not captured by the periodic model. Characterizing this transition, and constructing finite designs that approximate periodic behavior near boundaries, is an open direction.

\textit{Fabrication constraints.}
The framework is purely topological and does not explicitly account for material thickness, stiffness, collision avoidance, or manufacturing constraints. As a result, not all combinatorially valid structures are physically realizable without further processing. Bridging the gap between abstract topology and physical realization is a critical step toward practical applications, and incorporating such constraints into the construction process will enable the direct translation of combinatorial designs into realizable structures.

\begin{figure}[tb!]
    \centering
    \begin{subfigure}[b]{0.32\columnwidth}
        \includegraphics[width=0.99\linewidth]{images/laves/laves.jpg}\\
        \includegraphics[width=0.99\linewidth]{images/laves/laves_face_cycles.jpg}
    \end{subfigure}
    \begin{subfigure}[b]{0.64\columnwidth}
        \includegraphics[width=0.99\linewidth]{images/laves/laves_1.jpg}
    \end{subfigure}
    \caption{The Laves graph and the local structure of its vertices and edges. Although the Laves graph does not admit a canonical embedding, it exhibits combinatorial features suggestive of face constructions. In this particular embedding, each edge is incident to 10 distinct length-10 loops, and each vertex is surrounded by 15 such loops. We extract these loops and use them to define face cycles. The resulting face cycles are shown in the bottom-left image, while one of the possible induced LK structures is illustrated in the right image.}
    \Description{Three panels show a Laves graph, extracted local face cycles, and one corresponding linked-knot structure.}
    \label{fig:laves_intro}
\end{figure}

\textit{Bioinspired realizations.} The current design-level scope is particularly well suited to certain classes of structures that already exist in nature, such as \emph{Laves graphs} (Figure~\ref{fig:laves_intro}), which arise as edge graphs of space-filling polyhedra and occur ubiquitously in crystalline and intermetallic materials~\cite{berry1953crystal,coxeter1955laves}. Such graphs capture the underlying connectivity of a wide range of crystalline, molecular, and architected systems, where periodicity and uniform local coordination play a central role~\cite{hyde1997language,weaire1999physics}. Exploring how H-GRS-based designs can be matched to, and inform the analysis of, naturally occurring entangled structures is a promising direction.

\textit{Hierarchical and multi-scale extensions.} The combinatorial nature of the framework suggests natural extensions to hierarchical and multi-scale structures. Embedding scaffolds within scaffolds may provide a pathway toward designing complex architectured materials with tunable mechanical properties across scales.

\textit{Applications to architectured materials and mechanisms.} The unified treatment of articulated systems and filament networks suggests applications in soft robotics, deployable structures, textile engineering, and mechanical metamaterials. Pursuing these applications will further clarify the relationship between topology, geometry, and function, and will surface additional constraints that should feed back into the design framework.

\medskip

Overall, this work establishes a foundation for a topology-driven approach to structural design. By decoupling combinatorial structure from geometric realization, it opens a broad and largely unexplored design space at the intersection of geometry, topology, and fabrication.


\section{Conclusion}


We introduced a unified framework for generating articulated and filament-based structures based on hypergraph rotation systems augmented with local twist assignments. This formulation provides a combinatorial scaffold that encodes connectivity independently of geometric realization, enabling the systematic construction of a wide range of structures, including knots, links, chainmail systems, and periodic woven and knitted materials.

A central insight of this work is that structural complexity can be controlled through local combinatorial rules. By varying the underlying scaffold and twist parameters, one can continuously navigate a design space spanning unicursal filament systems, multi-cycle interlaced configurations, and hybrid structures. This establishes a direct correspondence between topological structure and classes of physical artifacts, bridging articulated mechanisms and architectured materials within a single paradigm.

Furthermore, the framework highlights the role of scaffold topology in determining global properties such as openness, closure, and periodicity. In particular, tree-based and cycle-based scaffolds lead to fundamentally different classes of articulated systems, while periodic scaffolds enable the construction of volumetric textile-like materials. This perspective provides a foundation for topology-driven design, where complex structures emerge from simple, local combinatorial specifications.

\begin{acks}
This work is partially supported by the National Science Foundation under Grant No. 2048182 and the J. Mike Walker Career Development Professorship received by Vinayak Krishnamurthy. Any opinions, findings, conclusions, or recommendations expressed in this material are those of the author(s) and do not necessarily reflect the views of the National Science Foundation or other funding sources.

\textit{AI Usage Statement.}
This work was developed through a collaboration between the authors and AI-assisted tools. AI was used to assist in drafting and refining the text. In one supplementary figure, the illustrations were initially created by an author, refined using AI tools, and subsequently corrected and finalized by the author. All conceptual contributions, theoretical results, and final decisions are solely those of the authors.
\end{acks}

\bibliographystyle{ACM-Reference-Format}
\bibliography{references}
\end{document}
